\chapter{Introducción}

\section{Introducción}

El agua es un recurso esencial para la vida en el planeta, y su calidad tiene un impacto directo en la salud de los ecosistemas, la biodiversidad y las actividades humanas. A nivel mundial, la contaminación hídrica ha alcanzado niveles críticos, comprometiendo el bienestar de millones de personas y el equilibrio de los ecosistemas acuáticos. En México, la situación es particularmente alarmante: menos del 25\% de las aguas residuales descargadas en cuerpos de agua reciben tratamiento adecuado, según el Fondo para la Comunicación y la Educación Ambiental, agravando significativamente el deterioro de la calidad del agua \cite{agua2006}.

En las últimas décadas, los métodos de monitoreo de la calidad del agua han avanzado, pero aún dependen en gran medida de técnicas manuales y análisis de laboratorio. En México, diversos programas buscan involucrar a las comunidades en la detección de cambios en la calidad del agua mediante mediciones rápidas y variables básicas, apoyándose en guías que estandarizan metodologías y validan resultados. Sin embargo, estas mediciones son limitadas; el proceso de recolección y validación es prolongado, y se lleva a cabo en periodos definidos, lo que impide un monitoreo constante.

Para realizar este tipo de monitoreos participativos, los interesados necesitan seguir una serie de pasos para poder ser aceptados por las instituciones responsables. Estos requerimientos incluyen presentar una solicitud formal, cumplir con los requisitos de capacitación, obtener validación técnica y comprometerse a seguir protocolos estrictos. Además, es necesario garantizar que los datos recopilados sean fiables y compatibles con los estándares definidos, lo que implica una evaluación continua por parte de capacitadores acreditados y la supervisión institucional. Si bien este enfoque busca fomentar la participación ciudadana y garantizar la utilidad de los resultados para la gestión ambiental, presenta desafíos importantes. La participación comunitaria tiende a disminuir con el tiempo debido a la falta de incentivos, recursos y reconocimiento. Sumado a esto, los datos recolectados pueden no ser aceptados oficialmente si carecen del rigo requerido; e incluso, para generar una ``primera alerta'' se exige un historial acumulado de monitoreo y evidencias de respaldo para corroborar el muestreo.

Una situación similar ocurre con los monitoreos realizados por las instituciones encargadas de garantizar el cumplimiento de estándares ambientales. Aunque generan reportes formales y detallados sobre el tipo y nivel de contaminantes, se trata de procesos lentos. Dependiendo del tipo de cuerpo de agua, la periodicidad del monitoreo varía (bimestral, semestral o anual). Consecuentemente, los resultados suelen retrasarse durante largos periodos antes de publicarse y, debido a su complejidad tecnica, resultan difíciles de comprender para el público general \cite{guiaMonitoreo2024}.

La falta de monitoreo constante en los cuerpos de agua que abastecen al país tiene un impacto severo en los ecosistemas y en la salud pública. Aunque las tecnologías actuales permiten medir parámetros clave como el pH, la oxigenación y la conductividad, estas mediciones continuan realizándose principalmente a través de muestreos manuales en lugar de métodos automatizados. Adicionalmente, contaminantes como los residuos sólidos no solo obstruyen los cuerpos de agua, sino que también actúan como vectores para la dispersión de contaminantes más peligrosos, como metales pesados y productos químicos.

Ante esta problematica la combinación de redes inalámbricas de sensores (WSN, por sus siglas en inglés) y visión artificial ofrece una solución prometedora para abordar las limitaciones de los métodos tradicionales. Las WSN permiten la automatización de la medición de parámetros clave como pH, oxigenación y conductividad de manera distribuida, lo que facilita un monitoreo más continuo y preciso del estado de la calidad del agua. Por su parte, la visión artificial mejora la detección de residuos sólidos, contribuyendo a identificar contaminantes visibles en los ríos. En conjunto, estos enfoques generan datos oportunos que pueden ayudar a las comunidades a comprender mejor las condiciones hídricas y fomentar acciones para su preservación.

Con base en esta propuesta, el sistema planteado (denominado ``MONICA'') busca integrar tecnologías avanzadas para automatizar y ofrecer un monitoreo constante de la calidad del agua en cuerpos lóticos\footnote{Lótico: Masa de agua con flujo constante en una dirección específica, como ríos, arroyos y corrientes\cite{conagua2024}.} (ríos, riachuelos y arroyos). A través de la recopilación de datos clave y su visualización en una plataforma accesible, se espera proporcionar información que facilite el análisis de las condiciones hídricas y la identificación de residuos flotantes. Este enfoque pretende ser un recurso valioso para quienes deseen vigilar el estado de los ecosistemas acuáticos, apoyando activamente a la difusión abierta de información ambiental.
